\chapter{Inégalités dans le triangle}

\begin{itemize}
\it On donne un point $M$ intérieur au triangle $ABC$. 
\begin{enumerate}
\item Montrer que $MA+MB$ est compris entre $AB$ et $CA+CB$.
\item En déduire que $MA+MB+MC$ est compris entre le demi-périmètre et le périmètre du triangle.
\end{enumerate}
\it Soit $O$ le point d'intersection des diagonales d'un quadrilatère convexe $ABCD$. \begin{enumerate}
\item Montrer que chaque diagonale est inférieure au demi-périmètre du quadrilatère.
\item Montrer que $AC+BD$ est supérieure à chacune des sommes $AB+CD$ et $AD+BC$.
\item En déduire que la somme des diagonales est comprise entre le demi-périmètre et le périmètre du quadrilatère.
\end{enumerate}
\it Démontrer que la médiane $AM$ d'un triangle $ABC$ est comprise entre la demi-différence et la demi-somme des côtés $AB$ et $AC$. (On prolongera $[AM]$ d'une longueur égale à elle-même.)
\it Soit un point $A$ et un cercle de centre $O$. Le diamètre passant par $A$ coupe le cercle entre $B$ et $C$. Soit $M$ un point quelconque du cercle. \begin{enumerate}
\item Comparer $AM$ à la somme et à la différence de $OA$ et $OM$.
\item Montrer que $AM$ est compris entre $AB$ et $AC$.
\end{enumerate}
\it Deux points $A$ et $B$ sont d'un même côté d'une droite $(xy)$. \begin{enumerate}
\item Trouver sur cette droite un point $P$, tel que la somme $PA+PB$, soit la plus petite possible (utiliser $A'$ symétrique de $A$ par rapport à $(xy)$.)
\item Que représente $(xy)$ pour l'angle $\widehat{APB}$ ?
\item Montrer que lorsque le point $M$ décrit la droite $(xy)$ la somme $MA+MB$ augmente en même temps que $PM$. 
\end{enumerate}
\it Deux points $A$ et $B$ sont de part et d'autre de $(xy)$. \begin{enumerate}
\item Trouver sur cette droite un point $P$, tel que la différence entre $PA$ et $PB$ soit la plus grande possible (utiliser $A'$ symétrique de $A$ par rapport à $(xy)$). 
\item Que représente $(xy)$ pour l'angle $\widehat{APB}$ ?
\end{enumerate}
\it On donne un angle $\widehat{xOy}$ inférieure à $60^o$ et deux points $A$ et $B$ intérieurs à cet angle. Trouver un point $M$ sur $(Ox)$ et un point $N$ sur $(Oy)$ de façon que le périmètre de la ligne brisée $AMNB$ soit le plus petit possible (utiliser les symétriques $A'$ de $A$ par rapport à $(Ox)$ et $B'$ de $B$ par rapport à $(Oy)$).
\it On considère un triangle $ABC$ et la bissectrice extérieure de l'angle en $A$. \begin{enumerate}
\item Montrer que le symétrique $C'$ de $C$ par rapport à cette bissectrice se trouve sur $(BA)$ et que $AC'=AC$. 
\item Démontrer que pour tout point $M$ de la bissectrice on a : \[ MB+MC>AB+AC.\]
\end{enumerate}
\it On considère un triangle $ABC$ et la bissectrice intérieure de l'angle en $A$. \begin{enumerate}
\item Montrer que le symétrique $C'$ de $C$ par rapport à cette bissectrice se trouve sur $(AB)$ et que $AC'=AC$. 
\item Démontrer que pour tout point $M$ de la bissectrice on a : \[ MB-MC>AB-AC.\]
\end{enumerate}
\it Dans un triangle $ABC$ on prolonge la médiane $[CM]$ d'une longueur $MD=CM$. \begin{enumerate}
\item Comparer les triangles $MBC$ et $MAD$ et montrer que $\widehat{BAD}=\widehat{ABC}$.
\item Soit $\widehat{BAx}$ l'angle extérieur en $A$ au
triangle $ABC$. Démontrer que $\widehat{BAx}>\widehat{ABC}$ et en déduire le théorème : \\
\emph{Dans tout triangle, un angle extérieur est supérieur à tout angle intérieur non adjacent.}
\end{enumerate} 
\it On considère deux triangles $ABC$ et $A'B'C'$ tels 
que $BC=B'C'$ et $\widehat{B}=\widehat{B'}$. On superpose ces deux triangles en amenant $B'$ en $B$, $C'$ en $C$ et $A'$ sur la demi-droite $[BA)$.\begin{enumerate}
\item En utilisant le théorème établi à l'exercice précédent, démontrer que l'une des inégalités $AB>A'B'$, $\widehat{C}>\widehat{C'}$ et $\widehat{A}<\widehat{A'}$ entraîne les deux autres.
\item Qu'arrive-t-il si $\widehat{A}=\widehat{A'}$ ? En déduire le cas d'égalité non classique : \\
\emph{Lorsque deux triangles ont deux angles homologues respectviement égaux et le côté opposé à l'un d'eux égal, ils sont égaux.}
\end{enumerate}
\it Soit un triangle $ABC$ dans lequel $AB<AC$. On prolonge la médiane $[AM]$ d'une longueur $MD=AM$.
\begin{enumerate}
\item Comparer les triangles $MCA$ et $MDB$. Conséquences ?
\item En déduire que : $\frac12(AC-AB)<AM<\frac12(AB+AC)$. Comparer la somme des médianes au périmètre du triangle.
\item Démontrer que l'angle $\widehat{MAC}$ est inférieur à l'angle $\widehat{MAB}$ et que la bissectrice intérieure de l'angle $\widehat{BAC}$ est située à l'intérieur de l'angle $\widehat{MAB}$.
\end{enumerate}
\end{itemize}
